% prism-latex demo document.
% Compiles with a stock TeX Live installation using pdflatex and produces
% several pages, which makes it useful for verifying the compile-and-rasterize
% pipeline end to end.
\documentclass[11pt]{article}

\usepackage[a4paper,margin=1in]{geometry}
\usepackage{amsmath}
\usepackage{amssymb}

\newcommand{\prism}{\textsc{Open-Prism}}
\newcommand{\dsh}{\textsc{DeepSeek Harness}}

\title{\prism{} in \dsh:\\ A Living \LaTeX{} Preview}
\author{The \texttt{prism-latex} bundle}
\date{Version 0.1.0}

\begin{document}

\maketitle

\begin{abstract}
This sample document exists so the preview can be checked against a real
compiler rather than a mock. It exercises the parts of a small paper that a
renderer is likely to meet: a title block, an abstract, numbered sections,
lists, displayed mathematics, a table, and enough prose to spill onto a second
page. If this file compiles and its pages appear in the right sidebar, the
pipeline from source to raster image is working.
\end{abstract}

\section{Why compile inside the sidebar?}

A language model writes \LaTeX{} the way a person does: as source. Reading that
source is useful, but it is not the same as reading the document. The difference
is not cosmetic. Page breaks move, a table overflows, an equation number
collides with a line of text, and a stray backslash turns a finished section
into an error message. None of those problems are visible in a syntax-highlighted
editor, and all of them are obvious in a compiled page.

The usual workaround is a terminal: run the engine, open the PDF, scroll to the
place that changed. That loop is fine on a workstation with a build script and
poor inside a chat window. The bundle this document tests closes the loop in
place: opening a \texttt{.tex} file compiles it with the TeX Live installation
already on the machine and paints the pages in the document preview.

\section{The pipeline}

The work is split between two halves that never share a process.

\begin{enumerate}
  \item The host half resolves the requested file, checks that it is a
        \texttt{.tex}, \texttt{.latex}, or \texttt{.ltx} file inside the
        session workspace, and refuses anything that escapes that directory.
  \item It creates a throwaway directory and runs the engine there, with shell
        escape disabled. The file's own directory is the working directory, so
        \verb|\input| and relative image paths resolve the way they would in a
        local build.
  \item It reads the produced PDF and hands it to Ghostscript, which rasterizes
        the first pages at the configured resolution.
  \item It returns the pages as PNG data with their pixel dimensions, plus a
        freshness token and a link to the PDF.
  \item The client half turns each page into a blob URL, stacks them vertically,
        and keeps the tab's loading state honest: compiling, failed, or ready.
\end{enumerate}

A small in-memory cache is keyed by the file's real path, its modification time,
its size, and the engine name. Reopening an unchanged document reuses the
previous compile; editing it invalidates the entry automatically.

\section{A little mathematics}

Because a preview is only convincing when the hard parts render correctly, the
document includes a few constructs that stress font and spacing decisions:
inline mathematics such as $e^{i\pi} + 1 = 0$, a display with an integral, and
an aligned system.

\begin{equation}
  \int_{-\infty}^{\infty} e^{-x^{2}}\,\mathrm{d}x = \sqrt{\pi}.
\end{equation}

\begin{align}
  \nabla \cdot \mathbf{E} &= \frac{\rho}{\varepsilon_{0}}, &
  \nabla \cdot \mathbf{B} &= 0, \\
  \nabla \times \mathbf{E} &= -\frac{\partial \mathbf{B}}{\partial t}, &
  \nabla \times \mathbf{B} &= \mu_{0}\mathbf{J}
    + \mu_{0}\varepsilon_{0}\frac{\partial \mathbf{E}}{\partial t}.
\end{align}

The Gaussian integral above is the only numbered result the sample needs. The
alignment is here to make sure that a two-column display does not push the text
off the page. It is written for a single compiler pass: the preview runs one
engine invocation, so the sample avoids cross-references and tables of contents
that would need a second pass to settle.

\section{What the preview shows}

The page stack is deliberately plain, so that the document, not the interface,
carries the attention. The table below lists the controls and what they do.

\begin{table}[h]
  \centering
  \begin{tabular}{ll}
    \hline
    Control & Effect \\
    \hline
    Fit width & Default; each page matches the pane width \\
    Zoom out / in & Fixed steps from 50\% to 300\% \\
    Percentage & Current scale, or \emph{Fit} while fitting \\
    Page count & Pages rendered, and the compile time \\
    Download PDF & Saves the compiled vector PDF \\
    Retry & Recompiles after an error \\
    \hline
  \end{tabular}
  \caption{Controls in the preview pill, in the order they appear.}
\end{table}

Fit width never enlarges a page beyond its natural size on screen. That matters
on a wide monitor: a short page should look like paper, not like a poster. Fixed
zoom, by contrast, is allowed to overflow the pane horizontally, and the pane
scrolls when it does.

\section{Limits worth knowing}

\begin{itemize}
  \item The preview is read-only. Editing happens in the source file; a new
        version of the file, or the pane's Reload button, compiles it again.
  \item Pages are raster images. Text is not selectable in the pane; the PDF
        button is the way to obtain the vector original.
  \item Shell escape is always disabled, so packages that need to run external
        programs will not compile.
  \item Long documents preview their first pages only. The PDF route still
        serves the complete file.
\end{itemize}

These are honest limits of a preview, not of \LaTeX{}. A compiled document is
still the best proof that the source is what its author meant, and a preview
that stays one click away from the source is enough to keep the loop short.

\section*{Acknowledgements}

The idea of pairing an AI writing workspace with a compiled document view comes
from the open-source \prism{} project. This bundle ports that idea into the
\dsh{} plugin system; the compiler and the rasterizer remain the ones already
installed on the machine.

\end{document}
